\scanpage{31}
\mspage{1}

Comparaison entre $\mathrm{Cat}$ et $\unsure{\Delta^{\wedge}|\widehat{\Delta}|\Delta^{1}}$

1. \deleted{Les notions de} \inserted{Rappel des notions de} $\mathbb D$-asphéricité\unsure{,|;} \deleted{dans $\Delta^{\wedge}$, pour $A$ une petite catégorie.} \inserted{et lien avec le ``\uline{localisateur fondamental}'' $W_{\mathbb D}$.}

Pour un dérivateur $\mathbb D_{\unsure{*|\bullet}}$, il s'agit des notions suivantes (cf I) :

(i) \uline{Objet} $\mathbb D$-\uline{asphérique} de $\mathrm{Cat}$, (i.e $Z\longrightarrow e$ est une $\mathbb D$-équivalence)

NB Je préfère dire ``$\mathbb D$-asphérique'' au lieu de ``$\mathbb D$-contractile'', qui prête parfois à confusion.

\textcircled{\scriptsize ii} $\mathbb D$-équivalence pour une flèche $f:X\mathbin{\unsure{-|\longrightarrow}}Y$ dans $\mathrm{Cat}$

(iii) \uline{Foncteurs} $f:X\longrightarrow Y$ $\mathbb D$-\uline{asphériques} (i.e. les $X_{/y}$ sont $\mathbb D$-asphériques) et $\mathbb D$-coasphériques (les $\,{}_{y\backslash}X$ sont $\mathbb D$-asphériques). C'est ce que j'avais appelé ``$\mathbb D$-équivalence stricte'' (ou ``costricte''), ou foncteur \inserted{\unsure{homologiquement|cohomologiquement}} $\mathbb D$-cofinal$^{*}$ (ou cohomologiquement $\mathbb D$-cofinal) \deleted{\unsure{pour}}, mais la terminologie ``asphérique'' me semble la plus simple et la meilleure.

(\unsure{Pêcher|Pécher|Pêché} dans PS, Pursuing Stacks)

(iv) \uline{Foncteurs} $f:X\longrightarrow Y$ $\mathbb D$-\uline{lisses}, $\mathbb D$-\uline{propres}, i.e. tels que les fibres des foncteurs induits $X_{/x}\longrightarrow Y_{/f(x)}$, respt. $\,{}_{x\backslash}X\longrightarrow{}_{f(x)\backslash}Y$, soient $\mathbb D$-asphériques

(v) $\mathbb D$-\uline{fibration} $f:X\longrightarrow Y$ (notion autoduale), i.e. telle que pour $Y''\xrightarrow{u}Y'$ au dessus de $Y$, si \inserted{$Y'',Y'$ \deleted{\unsure{les}} \inserted{(Hot)} asph.} \deleted{$u$ est une $\mathbb D$-\unsure{équi}}, alors \deleted{il en est de même de} $u_{X'}:X''\longrightarrow X'$ dans

\diagram{\begin{tikzcd}[ampersand replacement=\&,column sep=small,row sep=small]
 X\arrow[d]\&X'\arrow[l]\arrow[d]\&X''\arrow[l]\arrow[d]\\
 Y\&Y'\arrow[l]\&Y''\arrow[l]
\end{tikzcd}}

est une $\mathbb D$-équiv.

% slanted
\margin{On dira aussi que $f$ est \uline{$\mathbb D$-fibrant}. (NB Il n'y a pas de notion ``$\mathbb D$-cofibrant'' à ma connaissance.)}

La $\mathbb D$-\uline{fibration} \deleted{\unsure{est}} \inserted{est} \uline{triviale} \inserted{(ou \uline{asphérique})} \unsure{sss|ssi} ses fibres sont asphériques, ou encore sss $f$ est asphérique, ou sss $f$ est coasphérique. Cela signifie aussi que $f$ est une $\mathbb D$-équiv. \uline{universelle}.
% archive sheet number: 30
